\section{Use word arithmetic inside full-width addresses}
Adding two algorithm words discards a carry out of bit 31. For example,
\texttt{ffffffff} plus one produces \texttt{00000000}. A signed-overflow exception
would interrupt a valid calculation, so use \texttt{addu}. Its result on MIPS64 is a
sign-extended word. The upper 32 register bits may therefore all be one even though
the algorithm stores only a 32-bit pattern. Compare the low word when observing the
calculation, and use \texttt{sw} to write its four bytes.

A left rotation moves the high bits that leave a word back into its low positions.
For a rotation by \(r\), where \(0<r<32\),
\[
  \operatorname{ROTL}_r(x)=
  ((x\ll r)\bmod2^{32})\mathbin{\mathrm{OR}}(x\gg(32-r)).
\]
The right shift here is logical. An arithmetic right shift can insert ones at the
left edge and corrupt the wraparound bits. MIPS supplies a word right rotate, so a
left rotate by 12 can be expressed as a right rotate by 20.
\Needspace{6\baselineskip}
\MARSCodeLabel{Rotate the low word left by twelve bits.}
\begin{minted}[samepage=true]{mars-idea}
rotr $t1, $t1, 20
\end{minted}

Alternatively, use two word shifts and combine their results. Both shifts read the
original word; storing the left shift back into \texttt{t1} before the right shift
would destroy the source needed for the second half.
\Needspace{8\baselineskip}
\MARSCodeLabel{Keep the original word until both shifts have read it.}
\begin{minted}[samepage=true]{mars-idea}
sll $t4, $t1, 12
srl $t5, $t1, 20
or $t1, $t4, $t5
\end{minted}

The final OR may leave a sign-extended upper half. The next word operation and the
four-byte store still use the intended low pattern. A doubleword rotate changes the
rotation width and is therefore unsuitable here.

Address calculations have a different rule. Increment a byte pointer with
\texttt{daddiu}, and combine a base with an offset using \texttt{daddu}. A word-sized
increment can truncate a valid address above 4 GiB. The fact that the pointed-to value
is a word does not narrow its address.
